% answers_short.tex —— 2002 年·春季上海卷 · 简短答案速查
% 由 tools/make_short_answers.py 自 answers.tex 的 \gksolhead/\gkanswer 提取，
% 并与 2002/2002_数学真题_春季上海_答案要点.md 逐题比对后生成，勿手改（改源 MD 后重跑）。

\begin{gknotice}
\textbf{使用说明}：本节只给最终答案便于快速核对，完整推导见解析版。答案为编写者独立推导，\textbf{非官方标准答案}。
\end{gknotice}

\gkpart{填空题}
\begin{itemize}[leftmargin=2.6em]
  \item[\textbf{1.}] $\mathbb R\setminus\{-3,1\}$。
  \item[\textbf{2.}] $\dfrac{(x-2)^2}{25}+\dfrac{y^2}{16}=1$。
  \item[\textbf{3.}] $P\cap Q^c$。
  \item[\textbf{4.}] $-1$。
  \item[\textbf{5.}] $18$。
  \item[\textbf{6.}] $\dfrac2{\sin\alpha}$。
  \item[\textbf{7.}] $\dfrac18$。
  \item[\textbf{8.}] $F_1(a,b)\ne0$ 或 $F_2(a,b)\ne0$。
  \item[\textbf{9.}] $\dfrac34$。
  \item[\textbf{10.}] $3$ 对。相互异面的线段对可记为 $EF$ 与 $HG$、$AB$ 与 $HG$、$AB$ 与 $CD$。
  \item[\textbf{11.}] $\dfrac{50\sqrt6}{3}\text{ \text{海里}}\approx40.8\text{ \text{海里}}$。
  \item[\textbf{12.}] $\displaystyle\frac{V_{OP_1Q_1R_1}}{V_{OP_2Q_2R_2}}=\frac{OP_1}{OP_2}\cdot\frac{OQ_1}{OQ_2}\cdot\frac{OR_1}{OR_2}$。
\end{itemize}

\gkpart{选择题}
\begin{center}\begin{tabular}{ccccc}
\toprule
题号 & 13 & 14 & 15 & 16 \\
\midrule
答案 & D & C & B & C \\
\bottomrule
\end{tabular}\end{center}

\gkpart{解答题}
\begin{itemize}[leftmargin=2.6em]
  \item[\textbf{17.}] $\omega=7-i$ 或 $\omega=-7+i$。
  \item[\textbf{18.}] $y=\pm\sqrt2x$。
  \item[\textbf{19.}] (1) $\arccos\dfrac1{2\sqrt2}$；(2) $\arccos\dfrac17$
  \item[\textbf{20.}] (1) 见解析；(2) 方程 $f(x)=0$ 没有负数根
  \item[\textbf{21.}] (1) $a_2=\dfrac{b(n-1)}{n^2}$，$a_3=\dfrac{b(n-1)^2}{n^3}$，$a_k=\dfrac bn\left(\dfrac{n-1}{n}\right)^{k-1}$；(2) $a_k>a_{k+1}$（$b>0$）；其实际意义是业绩排名越靠前，奖金越多；(3) $\displaystyle\lim_{n\to\infty}P_n(b)=\frac be$
  \item[\textbf{22.}] (1) 不动点为 $x=-1,3$；(2) $0<a<1$；(3) $b_{\min}=-\dfrac1{2\sqrt2}$
\end{itemize}
